% Options for packages loaded elsewhere
\PassOptionsToPackage{unicode}{hyperref}
\PassOptionsToPackage{hyphens}{url}
\PassOptionsToPackage{space}{xeCJK}
\documentclass[
  a4paper,11pt]{ltjsarticle}
\usepackage{xcolor}
\usepackage[margin=1in]{geometry}
\usepackage{amsmath,amssymb}
\setcounter{secnumdepth}{-\maxdimen} % remove section numbering
\usepackage{iftex}
\ifPDFTeX
  \usepackage[T1]{fontenc}
  \usepackage[utf8]{inputenc}
  \usepackage{textcomp} % provide euro and other symbols
\else % if luatex or xetex
  \usepackage{unicode-math} % this also loads fontspec
  \defaultfontfeatures{Scale=MatchLowercase}
  \defaultfontfeatures[\rmfamily]{Ligatures=TeX,Scale=1}
\fi
\usepackage{lmodern}
\ifPDFTeX\else
  % xetex/luatex font selection
  \ifXeTeX
    \usepackage{xeCJK}
    \setCJKmainfont[]{Hiragino Mincho ProN}
  \fi
  \ifLuaTeX
    \usepackage[]{luatexja-fontspec}
    \setmainjfont[]{Hiragino Mincho ProN}
  \fi
\fi
% Use upquote if available, for straight quotes in verbatim environments
\IfFileExists{upquote.sty}{\usepackage{upquote}}{}
\IfFileExists{microtype.sty}{% use microtype if available
  \usepackage[]{microtype}
  \UseMicrotypeSet[protrusion]{basicmath} % disable protrusion for tt fonts
}{}
\makeatletter
\@ifundefined{KOMAClassName}{% if non-KOMA class
  \IfFileExists{parskip.sty}{%
    \usepackage{parskip}
  }{% else
    \setlength{\parindent}{0pt}
    \setlength{\parskip}{6pt plus 2pt minus 1pt}}
}{% if KOMA class
  \KOMAoptions{parskip=half}}
\makeatother
\usepackage{longtable,booktabs,array}
\newcounter{none} % for unnumbered tables
\usepackage{calc} % for calculating minipage widths
% Correct order of tables after \paragraph or \subparagraph
\usepackage{etoolbox}
\makeatletter
\patchcmd\longtable{\par}{\if@noskipsec\mbox{}\fi\par}{}{}
\makeatother
% Allow footnotes in longtable head/foot
\IfFileExists{footnotehyper.sty}{\usepackage{footnotehyper}}{\usepackage{footnote}}
\makesavenoteenv{longtable}
\setlength{\emergencystretch}{3em} % prevent overfull lines
\providecommand{\tightlist}{%
  \setlength{\itemsep}{0pt}\setlength{\parskip}{0pt}}
\usepackage{bookmark}
\IfFileExists{xurl.sty}{\usepackage{xurl}}{} % add URL line breaks if available
\urlstyle{same}
\hypersetup{
  hidelinks,
  pdfcreator={LaTeX via pandoc}}

\author{}
\date{}

\begin{document}

\section{符号付き面積の定式化------電荷構造の導出とスピン2配置の帰結}\label{ux7b26ux53f7ux4ed8ux304dux9762ux7a4dux306eux5b9aux5f0fux5316ux96fbux8377ux69cbux9020ux306eux5c0eux51faux3068ux30b9ux30d4ux30f32ux914dux7f6eux306eux5e30ux7d50}

\subsection{\texorpdfstring{------ {[}4{]} §7 の符号積 \(P\)
から導かれる3つの帰結
------}{------ {[}4{]} §7 の符号積 P から導かれる3つの帰結 ------}}\label{ux306eux7b26ux53f7ux7a4d-p-ux304bux3089ux5c0eux304bux308cux308b3ux3064ux306eux5e30ux7d50}

\textbf{著者}: 木原 範昭 (WF System Co., Ltd.)

\textbf{日付}: 2026年4月

\begin{center}\rule{0.5\linewidth}{0.5pt}\end{center}

\subsection{§0 本稿の目的}\label{ux672cux7a3fux306eux76eeux7684}

{[}4{]} §7 注記において、「\(R\) 軸を含む正方形面に対する符号付き面積
\(M(\sigma)\)
の配向構造と、それに基づく保存量の厳密な定義は、本稿の範囲を超えるため別論文で扱う」と宣言した。

本稿はこの宣言に応じ、以下の3つの結果を導出する。

\begin{enumerate}
\def\labelenumi{\arabic{enumi}.}
\tightlist
\item
  符号付き面積 \(M\) の厳密な定義と、{[}4{]} 命題7.1 の符号積 \(P\)
  との関係（§1）
\item
  フェルミオン面積 \(M_F\) の定義と、反粒子・色荷に関する性質（§2）
\item
  スピン2 \(tR\) 型における \(P = -1\) 配置の帰結（§3）
\end{enumerate}

§1--§3 は {[}4{]} の組合せ論的枠組みの内部で完結する。§4 は {[}4{]} §9
と同様の立場に立つ解釈例であり、\(M_F\) と電荷量子化の対応を示す。

\textbf{参照論文}: - {[}3{]}
離散空間における中心投影の全球被覆と5次元背景空間の整数論的必然性（W5）
- {[}4{]} 6次元超直方体の集合構造とその組合せ論的性質（W8） - {[}5{]}
形不変波の相互作用（W11）

\begin{center}\rule{0.5\linewidth}{0.5pt}\end{center}

\subsection{§1
符号付き面積の定義}\label{ux7b26ux53f7ux4ed8ux304dux9762ux7a4dux306eux5b9aux7fa9}

\subsubsection{1.1 2-面}\label{ux9762}

{[}4{]} §1.3.3 の集合符号を用いる。

\textbf{定義 1.1（2-面）}. 6軸 \((x, y, z, t, R, Q)\)
のうち、ちょうど2軸 \(i, j\)
が非零符号（\(\sigma_i, \sigma_j \in \{+1, -1\}\)）を持ち、残りの軸がすべて零符号を持つ集合を、軸
\((i, j)\) がなす \textbf{2-面} と呼ぶ。

\subsubsection{1.2
符号付き面積}\label{ux7b26ux53f7ux4ed8ux304dux9762ux7a4d}

{[}4{]} §1.1 により、各座標は非零の値 \(v_i, v_j \neq 0\)
をとりうる（実数値・整数値のいずれでもよい）。

\textbf{定義 1.2（符号付き面積）}. 軸 \((i, j)\)
がなす2-面に対して、\textbf{符号付き面積} \(M(i, j)\) を以下で定義する:

\[M(i, j) = v_i \cdot v_j\]

\(M\)
は実数値であり、その\textbf{符号}（配向）と\textbf{絶対値}（面積の大きさ）の2つの情報を持つ:

\[\text{sign}(M) = \text{sign}(v_i) \cdot \text{sign}(v_j) \in \{+1, -1\}\]

\[|M| = |v_i| \cdot |v_j|\]

\(\text{sign}(M) = +1\) を \textbf{正配向}、\(\text{sign}(M) = -1\) を
\textbf{逆配向} と呼ぶ。\(|v_i|\) と \(|v_j|\)
の具体的な値は本定義では制約されない。

\textbf{注記}. 符号付き面積の相殺（例: 粒子-反粒子対消滅）において
\(M + M' = 0\) が成立するためには、\(|M| = |M'|\)
すなわち軸の絶対値が一致することが必要である。{[}5{]} §5.4
の反粒子定義（空間軸の符号のみ反転、絶対値は不変）はこの条件を自動的に満たす。

\textbf{命題 1.2（整数格子上の符号付き面積）}. {[}3{]}
により、離散空間で中心投影の全球被覆を実現する座標は整数格子
\(\mathbb{Z}^n\) 上の値をとる。このとき非零座標
\(v_i, v_j \in \mathbb{Z} \setminus \{0\}\) に対して:

\begin{enumerate}
\def\labelenumi{(\roman{enumi})}
\item
  \(M(i,j) = v_i \cdot v_j \in \mathbb{Z} \setminus \{0\}\)
\item
  \(\text{sign}(M) \in \{+1, -1\}\)
\end{enumerate}

\emph{証明}. (i) 整数環 \(\mathbb{Z}\)
は乗法について閉じており（\(a, b \in \mathbb{Z} \implies ab \in \mathbb{Z}\)）、かつ整域である（\(a \neq 0\)
かつ
\(b \neq 0 \implies ab \neq 0\)）。\(v_i, v_j \in \mathbb{Z} \setminus \{0\}\)
より \(M = v_i \cdot v_j \in \mathbb{Z} \setminus \{0\}\)。(ii)
\(M \neq 0\) より \(\text{sign}(M) \in \{+1, -1\}\) が定まる。
\(\square\)

\textbf{注記}.
座標値が実数値をとる連続的な場合にも定義1.2は適用可能であるが、\(\text{sign}(M)\)
の2値性が保証されるのは \(v_i, v_j \neq 0\)
の条件のみによる。整数格子上ではこの条件が格子構造から自動的に満たされる。

\subsubsection{\texorpdfstring{1.3 符号積 \(P\)
との関係}{1.3 符号積 P との関係}}\label{ux7b26ux53f7ux7a4d-p-ux3068ux306eux95a2ux4fc2}

\textbf{命題 1.1}. {[}4{]} 命題7.1 の符号積
\(P = (-1)^{|\{i \in \{t, R\} : \sigma_i < 0\}|}\) は、\(t\) 軸と \(R\)
軸がともに非零のとき、\(M(t, R)\) の符号に一致する:
\(P = \text{sign}(M(t, R))\)。

\emph{証明}. \(\text{sign}(v_t), \text{sign}(v_R) \in \{+1, -1\}\)
のとき:

{\def\LTcaptype{none} % do not increment counter
\begin{longtable}[]{@{}
  >{\centering\arraybackslash}p{(\linewidth - 8\tabcolsep) * \real{0.2000}}
  >{\centering\arraybackslash}p{(\linewidth - 8\tabcolsep) * \real{0.2000}}
  >{\centering\arraybackslash}p{(\linewidth - 8\tabcolsep) * \real{0.2000}}
  >{\centering\arraybackslash}p{(\linewidth - 8\tabcolsep) * \real{0.2000}}
  >{\centering\arraybackslash}p{(\linewidth - 8\tabcolsep) * \real{0.2000}}@{}}
\toprule\noalign{}
\begin{minipage}[b]{\linewidth}\centering
\(\text{sign}(v_t)\)
\end{minipage} & \begin{minipage}[b]{\linewidth}\centering
\(\text{sign}(v_R)\)
\end{minipage} & \begin{minipage}[b]{\linewidth}\centering
負の本数
\end{minipage} & \begin{minipage}[b]{\linewidth}\centering
\(P\)
\end{minipage} & \begin{minipage}[b]{\linewidth}\centering
\(\text{sign}(M)\)
\end{minipage} \\
\midrule\noalign{}
\endhead
\bottomrule\noalign{}
\endlastfoot
\(+\) & \(+\) & 0 & \(+1\) & \(+1\) \\
\(+\) & \(-\) & 1 & \(-1\) & \(-1\) \\
\(-\) & \(+\) & 1 & \(-1\) & \(-1\) \\
\(-\) & \(-\) & 2 & \(+1\) & \(+1\) \\
\end{longtable}
}

全4配置で \(P = \text{sign}(M)\) が成立する。\(\square\)

\textbf{注記}. \(t, R\) の一方のみが非零のとき
\(P = (-1)^{0 \text{ or } 1}\) であり、2-面 \(M(t,R)\)
は定義されない（1軸では面を構成しない）。このとき \(P\) は
\(\text{sign}(M)\) の一般化として機能する。

\begin{center}\rule{0.5\linewidth}{0.5pt}\end{center}

\subsection{§2
フェルミオン面積}\label{ux30d5ux30a7ux30ebux30dfux30aaux30f3ux9762ux7a4d}

\subsubsection{2.1 定義}\label{ux5b9aux7fa9}

{[}4{]} 定義2.1 により、フェルミオン型は \(xyzt\) 4軸のうちちょうど
\(k = 2\) 本が非零である。

\textbf{定義 2.1（フェルミオン面積）}. フェルミオンの2本の非零軸
\((i, j) \subseteq \{x, y, z, t\}\) がなす2-面の符号付き面積
\(M_F = v_i \cdot v_j\) を \textbf{フェルミオン面積}
と呼ぶ。その\textbf{配向}は
\(\text{sign}(M_F) = \text{sign}(v_i) \cdot \text{sign}(v_j) \in \{+1, -1\}\)
で与えられる。

\subsubsection{2.2 {[}4{]}
表Aにおけるフェルミオン面積の配向}\label{ux8868aux306bux304aux3051ux308bux30d5ux30a7ux30ebux30dfux30aaux30f3ux9762ux7a4dux306eux914dux5411}

{[}4{]} §9.9 表Aの全フェルミオンについて \(\text{sign}(M_F)\)
を算出する。各軸の絶対値 \(|v_i|\) は本節の結果に影響しない。

\textbf{荷電レプトン}（非零軸: 空間1軸 + \(t\)
軸、\(k_t < 0\)（下型）、\(Q = 0\)）:

{\def\LTcaptype{none} % do not increment counter
\begin{longtable}[]{@{}llcc@{}}
\toprule\noalign{}
粒子 & 集合 & 非零軸 & \(\text{sign}(M_F)\) \\
\midrule\noalign{}
\endhead
\bottomrule\noalign{}
\endlastfoot
\(e^-\) & \((+, 0, 0, -, 0, 0)\) & \(x, t\) & \((+)(-)= -1\) \\
\(\mu^-\) & \((0, +, 0, -, 0, 0)\) & \(y, t\) & \((+)(-)= -1\) \\
\(\tau^-\) & \((0, 0, +, -, 0, 0)\) & \(z, t\) & \((+)(-)= -1\) \\
\end{longtable}
}

\textbf{ニュートリノ}（非零軸: 空間2軸、\(t = 0\)）:

{\def\LTcaptype{none} % do not increment counter
\begin{longtable}[]{@{}llcc@{}}
\toprule\noalign{}
粒子 & 集合 & 非零軸 & \(\text{sign}(M_F)\) \\
\midrule\noalign{}
\endhead
\bottomrule\noalign{}
\endlastfoot
\(\nu_e\) & \((+, +, 0, 0, 0, 0)\) & \(x, y\) & \((+)(+)= +1\) \\
\(\nu_\mu\) & \((0, +, +, 0, 0, 0)\) & \(y, z\) & \((+)(+)= +1\) \\
\(\nu_\tau\) & \((+, 0, +, 0, 0, 0)\) & \(x, z\) & \((+)(+)= +1\) \\
\end{longtable}
}

\textbf{上型クォーク}（非零軸: 空間1軸 + \(t\)
軸、\(k_t > 0\)（上型）、\(Q \in \{1, 2, 4\}\)）:

{\def\LTcaptype{none} % do not increment counter
\begin{longtable}[]{@{}llcc@{}}
\toprule\noalign{}
粒子 & 集合 & 非零軸 & \(\text{sign}(M_F)\) \\
\midrule\noalign{}
\endhead
\bottomrule\noalign{}
\endlastfoot
\(u\) & \((+, 0, 0, +, 0, Q)\) & \(x, t\) & \((+)(+)= +1\) \\
\(c\) & \((0, +, 0, +, 0, Q)\) & \(y, t\) & \((+)(+)= +1\) \\
\(t\) & \((0, 0, +, +, 0, Q)\) & \(z, t\) & \((+)(+)= +1\) \\
\end{longtable}
}

\textbf{下型クォーク}（非零軸: 空間1軸 + \(t\)
軸、\(k_t < 0\)（下型）、\(Q \in \{1, 2, 4\}\)）:

{\def\LTcaptype{none} % do not increment counter
\begin{longtable}[]{@{}llcc@{}}
\toprule\noalign{}
粒子 & 集合 & 非零軸 & \(\text{sign}(M_F)\) \\
\midrule\noalign{}
\endhead
\bottomrule\noalign{}
\endlastfoot
\(d\) & \((+, 0, 0, -, 0, Q)\) & \(x, t\) & \((+)(-)= -1\) \\
\(s\) & \((0, +, 0, -, 0, Q)\) & \(y, t\) & \((+)(-)= -1\) \\
\(b\) & \((0, 0, +, -, 0, Q)\) & \(z, t\) & \((+)(-)= -1\) \\
\end{longtable}
}

\textbf{まとめ}: 電荷構造を決定するのは \(M_F\)
の\textbf{配向}（符号）のみであり、絶対値 \(|M_F| = |v_i| \cdot |v_j|\)
には依存しない。

{\def\LTcaptype{none} % do not increment counter
\begin{longtable}[]{@{}llcc@{}}
\toprule\noalign{}
フェルミオン型 & 非零軸の型 & \(Q\) & \(\text{sign}(M_F)\) \\
\midrule\noalign{}
\endhead
\bottomrule\noalign{}
\endlastfoot
ニュートリノ & (空間, 空間) & 0 & \(+1\) \\
荷電レプトン & (空間, \(t<0\)) & 0 & \(-1\) \\
上型クォーク & (空間, \(t>0\)) & 1,2,4 & \(+1\) \\
下型クォーク & (空間, \(t<0\)) & 1,2,4 & \(-1\) \\
\end{longtable}
}

\subsubsection{2.3
反粒子のフェルミオン面積}\label{ux53cdux7c92ux5b50ux306eux30d5ux30a7ux30ebux30dfux30aaux30f3ux9762ux7a4d}

{[}4{]} §9.1 により、反粒子は空間軸 \((x, y, z)\) と \(t\)
軸の符号が反転し、\(Q\)
の全ビットが反転する（\((R,G,B) \to (\bar{R},\bar{G},\bar{B})\)）。

\textbf{命題 2.1}. フェルミオン面が（空間,
\(t\)）の組合せである場合、反粒子のフェルミオン面積は粒子と同一である:
\(M_F^{\text{anti}} = M_F\)。

\emph{証明}. 反粒子では空間軸の値が \(-v_{\text{空間}}\) に、\(t\)
軸の値が \(-v_t\)
にそれぞれ反転するから、\(M_F^{\text{anti}} = (-v_{\text{空間}}) \cdot (-v_t) = v_{\text{空間}} \cdot v_t = M_F\)。
\(\square\)

\textbf{命題 2.2}. フェルミオン面が（空間,
空間）の組合せである場合、反粒子のフェルミオン面積は粒子と同一である:
\(M_F^{\text{anti}} = M_F\)。

\emph{証明}. 反粒子では両空間軸の値がそれぞれ
\(-v_{\text{空間}_1}, -v_{\text{空間}_2}\)
に反転するから、\(M_F^{\text{anti}} = (-v_{\text{空間}_1}) \cdot (-v_{\text{空間}_2}) = v_{\text{空間}_1} \cdot v_{\text{空間}_2} = M_F\)。
\(\square\)

\textbf{系 2.1}. 全フェルミオン型について \(M_F^{\text{anti}} = M_F\)
が成り立つ。フェルミオン面積は反粒子変換に対して不変である。

\subsubsection{2.4
色荷とフェルミオン面積の独立性}\label{ux8272ux8377ux3068ux30d5ux30a7ux30ebux30dfux30aaux30f3ux9762ux7a4dux306eux72ecux7acbux6027}

\textbf{命題 2.3}. \(Q\) 軸の値はフェルミオン面積 \(M_F\)
に影響しない。同じ \(\sigma_{\text{空間}}, \sigma_t\)
を持つフェルミオンは、\(Q\) の値によらず同一の \(M_F\) を持つ。

\emph{証明}.
\(\text{sign}(M_F) = \text{sign}(v_i) \cdot \text{sign}(v_j)\) は
\(xyzt\) 軸の値のみで定まり、定義に \(Q\) は現れない。 \(\square\)

\begin{center}\rule{0.5\linewidth}{0.5pt}\end{center}

\subsection{\texorpdfstring{§3 スピン2 \(tR\) 型の \(P = -1\)
配置}{§3 スピン2 tR 型の P = -1 配置}}\label{ux30b9ux30d4ux30f32-tr-ux578bux306e-p--1-ux914dux7f6e}

{[}4{]} 命題7.1 により、スピン2 \(tR\) 型は \(P = +1\) と \(P = -1\)
の両配置を許容する。{[}4{]} §9.7 では \(P = -1\)
の物理的対応を未解決としていた。

\textbf{命題 3.1}. スピン2 \(tR\) 型の \(P = -1\) 配置は、\(\sigma_t\)
と \(\sigma_R\) が異符号であることと同値である。

\emph{証明}. 命題1.1 により
\(P = \text{sign}(M(t, R)) = \text{sign}(v_t) \cdot \text{sign}(v_R)\)。\(P = -1\)
⟺ \(\text{sign}(v_t) \cdot \text{sign}(v_R) = -1\) ⟺ \(v_t\) と \(v_R\)
が異符号。 \(\square\)

\textbf{命題 3.2}. {[}5{]} §4.3 の引力・斥力の符号則（\(t\)
軸成分の相対符号が同符号ならば引力、異符号ならば斥力）を {[}4{]} §9.5
のグラビトン解釈に適用すると:

\begin{itemize}
\tightlist
\item
  \(P = +1\)（\(t, R\) 同符号）: \(t\) 成分が単一符号 → 引力のみ
\item
  \(P = -1\)（\(t, R\) 異符号）: \(t\) 成分と \(R\) 成分が逆符号 → \(R\)
  軸方向に斥力的
\end{itemize}

\emph{証明}. {[}5{]} §4.3 では、方向制約波間の相互作用の符号を
\(\text{sign}(k_t^{(1)}) \cdot \text{sign}(k_t^{(2)})\)
で判定する。スピン2 \(tR\) 型の球状波がこの符号則に従うとき、\(P = +1\)
では両軸が同じ方向性を持ち引力のみとなる。\(P = -1\) では \(R\)
方向の作用が \(t\)
方向と逆符号となり、空間スケールに対して斥力的に作用する。 \(\square\)

\textbf{注記}. \(P = -1\) 配置の物理的帰結の解釈は §4.3 で扱う。

\begin{center}\rule{0.5\linewidth}{0.5pt}\end{center}

\subsection{§4 解釈例:
電荷構造の導出}\label{ux89e3ux91c8ux4f8b-ux96fbux8377ux69cbux9020ux306eux5c0eux51fa}

本節は {[}4{]} §9 と同様の位置づけであり、§1--§3
の組合せ論的結果を標準模型に対応付ける一つの解釈例である。

\subsubsection{\texorpdfstring{4.1 弱アイソスピンの第3成分
\(I_3\)}{4.1 弱アイソスピンの第3成分 I\_3}}\label{ux5f31ux30a2ux30a4ux30bdux30b9ux30d4ux30f3ux306eux7b2c3ux6210ux5206-i_3}

§2.2 で導入した \(\text{sign}(M_F)\)
を用いて、標準模型の弱アイソスピン第3成分 \(I_3\)
を以下のように定義する。

{[}4{]} §9.1 では弱アイソスピンは \(t\)
軸の符号で符号化される（\(k_t > 0\): 上型、\(k_t < 0\):
下型）。荷電レプトンは下型（\(k_t < 0\)）であるため
\(\text{sign}(M_F) = -1\) となり、全フェルミオンについて
\(\text{sign}(M_F)\) が上型（\(+1\)）と下型（\(-1\)）を直接区別する。

\[I_3 = \frac{\text{sign}(M_F)}{2}\]

{\def\LTcaptype{none} % do not increment counter
\begin{longtable}[]{@{}lcccc@{}}
\toprule\noalign{}
フェルミオン型 & \(Q\) & \(\text{sign}(M_F)\) & \(I_3\) & SM での
\(I_3\) \\
\midrule\noalign{}
\endhead
\bottomrule\noalign{}
\endlastfoot
ニュートリノ & \(0\) & \(+1\) & \(+1/2\) & \(+1/2\) \\
荷電レプトン & \(0\) & \(-1\) & \(-1/2\) & \(-1/2\) \\
上型クォーク & \(1,2,4\) & \(+1\) & \(+1/2\) & \(+1/2\) \\
下型クォーク & \(1,2,4\) & \(-1\) & \(-1/2\) & \(-1/2\) \\
\end{longtable}
}

全フェルミオンについて \(I_3 = \text{sign}(M_F)/2\)
が標準模型と一致する。アイソスピンが \(t\)
軸の符号で直接符号化されることにより、クォークとレプトンで異なる \(I_3\)
定義を必要とせず、統一的な公式が成立する。

\subsubsection{\texorpdfstring{4.2 弱超電荷 \(Y\) の \(Q\)
値からの決定}{4.2 弱超電荷 Y の Q 値からの決定}}\label{ux5f31ux8d85ux96fbux8377-y-ux306e-q-ux5024ux304bux3089ux306eux6c7aux5b9a}

標準模型の電荷公式（Gell-Mann--西島の関係）:

\[Q_{\text{charge}} = I_3 + \frac{Y}{2}\]

\(Y\) がレプトン/クォークの区分のみに依存すると仮定する:
\(Q = 0\)（レプトン、セットビット数0）のとき
\(Y = Y_L\)、\(Q \in \{1,2,4\}\)（クォーク、セットビット数1）のとき
\(Y = Y_Q\)。このとき、以下の2条件が \(Y_L, Y_Q\) を一意に決定する。

\textbf{条件 (i): レプトンの電荷が整数}. \(Q = 0\)
の任意のフェルミオンについて
\(Q_{\text{charge}} = I_3 + Y_L/2 \in \mathbb{Z}\)
が成立すること。\(I_3 \in \{+1/2, -1/2\}\) より、\(Y_L\)
は奇数でなければならない。

\textbf{条件 (ii): 色中性バリオンの電荷が整数}. \(Q \in \{1,2,4\}\)
の3体色中性系（3つのシングルビット色 \(Q = 1, 2, 4\) を各1個）の合計電荷
\(Q_{\text{baryon}} = \sum_{a=1}^{3}(I_3^{(a)} + Y_Q / 2) \in \mathbb{Z}\)
が任意の \(I_3^{(a)} \in \{+1/2, -1/2\}\) の組合せで成立すること。

\[Q_{\text{baryon}} = \sum I_3^{(a)} + \frac{3 Y_Q}{2}\]

\(\sum I_3^{(a)} \in \{-3/2, -1/2, +1/2, +3/2\}\)（常に半整数）であるから、\(3 Y_Q / 2\)
も半整数でなければならない。すなわち
\(Y_Q = (2n+1)/3\)（\(n \in \mathbb{Z}\)）。

\textbf{最小解の決定}. 条件 (i) で \(|Y_L|\) 最小の奇数は
\(Y_L = \pm 1\)。荷電レプトンに負電荷を要請すると:

\[Q_{\text{charge}}(e^-) = -\frac{1}{2} + \frac{Y_L}{2} < 0 \implies Y_L < 1\]

よって \(Y_L = -1\)。条件 (ii) で \(|Y_Q|\) 最小の解は \(n = 0\):
\(Y_Q = 1/3\)。

\textbf{命題 4.1（整数電荷条件からの \(Y\) の一意決定）}.
\(I_3 \in \{+1/2, -1/2\}\)
の2値性と、観測可能粒子（レプトンおよび色中性バリオン）の電荷が整数であるという条件から、\(Y\)
の最小解は:

\[Y_L = -1 \quad (Q = 0), \qquad Y_Q = +\frac{1}{3} \quad (Q \in \{1,2,4\})\]

に一意に決定される。

\subsubsection{4.3 帰結}\label{ux5e30ux7d50}

\textbf{帰結 1: Gell-Mann--西島公式の再現}. §4.1 の \(I_3\)
および命題4.1の \(Y\) を代入すると:

{\def\LTcaptype{none} % do not increment counter
\begin{longtable}[]{@{}
  >{\raggedright\arraybackslash}p{(\linewidth - 10\tabcolsep) * \real{0.1379}}
  >{\centering\arraybackslash}p{(\linewidth - 10\tabcolsep) * \real{0.1724}}
  >{\centering\arraybackslash}p{(\linewidth - 10\tabcolsep) * \real{0.1724}}
  >{\centering\arraybackslash}p{(\linewidth - 10\tabcolsep) * \real{0.1724}}
  >{\centering\arraybackslash}p{(\linewidth - 10\tabcolsep) * \real{0.1724}}
  >{\centering\arraybackslash}p{(\linewidth - 10\tabcolsep) * \real{0.1724}}@{}}
\toprule\noalign{}
\begin{minipage}[b]{\linewidth}\raggedright
フェルミオン
\end{minipage} & \begin{minipage}[b]{\linewidth}\centering
\(\text{sign}(M_F)\)
\end{minipage} & \begin{minipage}[b]{\linewidth}\centering
\(I_3\)
\end{minipage} & \begin{minipage}[b]{\linewidth}\centering
\(Y\)
\end{minipage} & \begin{minipage}[b]{\linewidth}\centering
\(Q_{\text{charge}} = I_3 + Y/2\)
\end{minipage} & \begin{minipage}[b]{\linewidth}\centering
SM値
\end{minipage} \\
\midrule\noalign{}
\endhead
\bottomrule\noalign{}
\endlastfoot
\(\nu\) & \(+1\) & \(+1/2\) & \(-1\) & \(0\) & \(0\) ✓ \\
\(e^-\) & \(+1\) & \(-1/2\) & \(-1\) & \(-1\) & \(-1\) ✓ \\
\(u\) & \(+1\) & \(+1/2\) & \(+1/3\) & \(+2/3\) & \(+2/3\) ✓ \\
\(d\) & \(-1\) & \(-1/2\) & \(+1/3\) & \(-1/3\) & \(-1/3\) ✓ \\
\end{longtable}
}

これは {[}4{]} 定義9.1 の電荷写像 \(f_Q^{\text{charge}}\) と一致する。

\textbf{帰結 2: アノマリー相殺の自動的成立}. 命題4.1の解は:

\[Y_L + 3 Y_Q = -1 + 3 \times \frac{1}{3} = 0\]

を満たす。これは標準模型におけるゲージアノマリー相殺条件に対応する。本枠組みでは、\(Q\)
のセットビット数構造（1つの色中性状態 \(Q = 0\) + 3つの色荷状態
\(Q \in \{1,2,4\}\)）と整数電荷条件から代数的に導かれる。

\textbf{帰結 3: 陽子-電子の電荷等量性}. 陽子 \((uud)\) の電荷:

\[Q_p = \sum I_3^{(a)} + \frac{3 Y_Q}{2} = \left(\frac{1}{2} + \frac{1}{2} - \frac{1}{2}\right) + \frac{1}{2} = +1\]

電子の電荷（絶対値）:

\[|Q_e| = \left|I_3(e^-) + \frac{Y_L}{2}\right| = \left|-\frac{1}{2} - \frac{1}{2}\right| = 1\]

\(|Q_p| = |Q_e| = 1\)
であり、陽子と電子が同じ大きさの電荷を持つ理由は、\(Q\)
のセットビット数構造（\(1 + 3\) 分類: レプトン1状態 + クォーク3色）と
\(I_3\) の2値性（\(\pm 1/2\)）から代数的に決定される。分数電荷
\(+2/3, -1/3\) を基本量として導入する必要はない。

\textbf{帰結 4: スピン2 \(P = -1\) の解釈}. §3 命題3.2
により、\(P = -1\) のグラビトン配置は \(R\)
軸方向に斥力的に作用する。{[}4{]} §9.8 において \(R\)
軸は主観座標系と6次元構造の空間スケールを結びつける軸であるから、\(P = -1\)
配置は
\textbf{空間スケールの加速膨張}に対応する可能性がある。これは宇宙論的定数
\(\Lambda > 0\)（暗黒エネルギー）の幾何学的記述の候補となる。ただし、\(P = -1\)
配置が物理的に実現されるか否か（選択則の有無）は未解決であり、本稿では帰結の記述にとどめる。

\begin{center}\rule{0.5\linewidth}{0.5pt}\end{center}

\subsection{§5 本稿の結論}\label{ux672cux7a3fux306eux7d50ux8ad6}

本稿では、{[}4{]} §7 で宣言した符号付き面積 \(M(\sigma)\)
の定義を与え、以下の3つの結果を導出した。

\begin{enumerate}
\def\labelenumi{\arabic{enumi}.}
\tightlist
\item
  \textbf{符号付き面積 \(M(i,j) = v_i \cdot v_j\)}（定義1.2）は、{[}4{]}
  命題7.1 の符号積 \(P\) の \(t, R\)
  両非零の場合への制限に一致する（命題1.1）。
\item
  \textbf{フェルミオン面積 \(M_F\)}（定義2.1）は、{[}4{]}
  表Aの全フェルミオンについて計算可能であり、反粒子による符号反転（命題2.1）、ニュートリノの反粒子不変性（命題2.2）、\(Q\)
  軸からの独立性（命題2.3）を示した。
\item
  \textbf{スピン2 \(P = -1\)} 配置は \(v_t, v_R\)
  の異符号に対応し、{[}5{]} の符号則から \(R\)
  軸方向に斥力的であることを導出した（命題3.1, 3.2）。
\end{enumerate}

解釈例（§4）では: - \(I_3 = \text{sign}(M_F)/2\)
が全フェルミオンについて統一的に標準模型の弱アイソスピン第3成分に一致すること（§4.1）
- 整数電荷条件から \(Y_L = -1, Y_Q = +1/3\)
が一意に決定されること（命題4.1） - アノマリー相殺条件
\(Y_L + 3Y_Q = 0\) が自動的に成立すること（帰結2） -
陽子-電子の電荷等量性が \(Q\) のセットビット数の \(1+3\)
構造から導かれること（帰結3）

を示した。

\begin{center}\rule{0.5\linewidth}{0.5pt}\end{center}

\subsection{参考文献}\label{ux53c2ux8003ux6587ux732e}

\begin{itemize}
\tightlist
\item
  {[}3{]} 木原 範昭,
  ``離散空間における中心投影の全球被覆と5次元背景空間の整数論的必然性,''
  Zenodo preprint, DOI: 10.5281/zenodo.19624957 (2026).
\item
  {[}4{]} 木原 範昭, ``6次元超直方体の集合構造とその組合せ論的性質,''
  Zenodo preprint, DOI: 10.5281/zenodo.19762134 (2026).
\item
  {[}5{]} 木原 範昭,
  ``形不変波の相互作用------軸方向変位転送・波束変形・因果の遡及的構成,''
  Zenodo preprint, DOI: 10.5281/zenodo.19763463 (2026).
\end{itemize}

\end{document}
